\documentclass{ximera}
%\usepackage{helvet}
%\renewcommand{\familydefault}{\sfdefault}
\author{Jeffrey Kuan}
\title{Math 1050 Homework 3}
\license{CC: 0}
\begin{document}

\begin{abstract}
   The third math 1050 Homework 
\end{abstract}

\maketitle

Where an answer is written as a sign box followed by a fraction, enter a minus sign in the first box if the result is negative, and leave that box empty if it is positive.

\begin{problem}

To multiply two fractions, you must first rewrite the fractions as equivalent fractions with a common denominator.

\begin{multipleChoice}
\choice[correct]{Never True}
\choice{Always True}
\choice{Sometimes True}
\end{multipleChoice}

\end{problem}

\begin{problem}

An irrational number is a real number.

\begin{multipleChoice}
\choice[correct]{Always True}
\choice{Never True}
\choice{Sometimes True}
\end{multipleChoice}
\end{problem}

\begin{problem}

To write a decimal as a percent, multiply the decimal by 1/100.

\begin{multipleChoice}
\choice{Always True}
\choice{Sometimes True}
\choice[correct]{Never True}
\end{multipleChoice}

\end{problem}

Instructions for solutions:

The answer box in front will either be ``-'' or blank (if it's positive). Enter the numerator and denominator as integers.
\begin{problem}
Simplify
\[
\frac{1}{2}\left( - \frac{3}{4} \right) = \answer{-}\frac{\answer{3}}{\answer{8}}
\]
\end{problem}

\begin{problem}
Simplify
\[
\left(\frac{1}{2} \right) \left( - \frac{3}{4} \right)  \left( - \frac{5}{8} \right)  = \frac{\answer{15}}{\answer{64}}
\]
\end{problem}

\begin{problem}
Simplify
\[
\left( - \frac{3}{5}\right) \left( - \frac{2}{9} \right) = \frac{\answer{2}}{\answer{15}}
\]
\end{problem}

\begin{problem}
Simplify
\[
\left(-\frac{2}{7} \right) \left( - \frac{5}{6} \right)  \left( - \frac{3}{4} \right)  = \answer{-}\frac{\answer{5}}{\answer{28}}
\]
\end{problem}



\begin{problem}
Simplify
\[
\frac{1}{8} \div \left(-\frac{5}{12}\right) = \answer{-} \frac{\answer{3}}{\answer{10}}
\]
\end{problem}

\begin{problem}
Simplify
\[
\left(-\frac{4}{9}\right) \div \left(-\frac{2}{3}\right) = \frac{\answer{2}}{\answer{3}}
\]
\end{problem}

\begin{problem}
Simplify
\[
-\frac{4}{15} \div \left(-\frac{8}{5}\right) = \frac{\answer{1}}{\answer{6}}
\]
\end{problem}

\begin{problem}
Simplify
\[
\frac{9}{10} \div \left(-\frac{3}{5}\right) = \answer{-} \frac{\answer{3}}{\answer{2}}
\]
\end{problem}


\begin{problem}
Simplify
\[
\frac{4}{15} - \left(-\frac{8}{5}\right) = \frac{\answer{28}}{\answer{15}}
\]
\end{problem}

\begin{problem}
Simplify
\[
\left(\frac{2}{7}\right) - \frac{3}{5} = \answer{-} \frac{\answer{11}}{\answer{35}}
\]
\end{problem}

\begin{problem}
Simplify
\[
-\left(\frac{1}{2}\right) - \left(-\frac{5}{6}\right) = \frac{\answer{1}}{\answer{3}}
\]
\end{problem}

\begin{problem}
Simplify
\[
\frac{4}{9} - \frac{5}{6} = \answer{-} \frac{\answer{7}}{\answer{18}}
\]
\end{problem}

\begin{problem}
Simplify
\[
\frac{3}{8}+ \frac{2}{3} - \frac{1}{4} = \frac{\answer{19}}{\answer{24}}
\]
\end{problem}

\begin{problem}
Simplify
\[
\frac{4}{9} + \frac{5}{6} - \frac{1}{3} = \frac{\answer{17}}{\answer{18}}
\]
\end{problem}

\begin{problem}
Simplify
\[
-\frac{2}{5} + \frac{3}{10} - \frac{1}{2} = \answer{-} \frac{\answer{3}}{\answer{5}}
\]
\end{problem}

\begin{problem}
   Evaluate
\[
-3.07 - (-2.81) + 17.1 = \answer{16.84}
\]
\end{problem}

\begin{problem}
   Evaluate
\[
-7.46 - 2.89 + 15.7 = \answer{5.35}
\]
\end{problem}

\begin{problem} Evaluate
\[
6.15 - (-3.92) + (-8.4) = \answer{1.67}
\]
\end{problem}

\begin{problem} Evaluate 
\[
-4.23 - (-5.67) + 12.8 = \answer{14.24}
\]
\end{problem}

\begin{problem}
Write as a fraction (simplified) and as a decimal
\[
75\% = \frac{\answer{3}}{\answer{4}} \text{ and } 75\% = \answer{0.75}
\]
\end{problem}

\begin{problem}
Write as a fraction (simplified) and as a decimal
\[
64\% = \frac{\answer{16}}{\answer{25}} \text{ and } 64\% = \answer{0.64}
\]
\end{problem}

\begin{problem}
Write as a fraction (simplified) and as a decimal
\[
19\% = \frac{\answer{19}}{\answer{100}} \text{ and } 19\% = \answer{0.19}
\]
\end{problem}

\begin{problem}
Write as a fraction (simplified) and as a decimal
\[
450\% = \frac{\answer{9}}{\answer{2}} \text{ and } 450\% = \answer{4.5}
\]
\end{problem}

\begin{problem}
Write as a fraction 
\[
11 \frac{1}{9}\% = \frac{\answer{1}}{\answer{9}}
\]
\end{problem}

\begin{problem}
Write as a fraction
\[
\frac{1}{2}\% = \frac{\answer{1}}{\answer{200}}
\]
\end{problem}

\begin{problem}
Write as a decimal
\[
7.3\% = \answer{0.073}
\]
\end{problem}

\begin{problem}
Write as a decimal
\[
9.15\% = \answer{0.0915}
\]
\end{problem}

\begin{problem}
Write as a percent
\[
0.15 = \answer{15\%}
\]
\end{problem}

\begin{problem}
Write as a percent
\[
0.175 = \answer{17.5\%}
\]
\end{problem}

\begin{problem}
Write as a percent
\[
0.008 = \answer{0.8\%}
\]
\end{problem}

\begin{problem}
Write as a percent. Round to the nearest tenth of a percent.
\[
\frac{4}{9} \approx \answer{44.4\%}
\]
\end{problem}

\begin{problem}
Write as a percent. Round to the nearest tenth of a percent.
\[
\frac{27}{50} = \answer{54\%}
\]
\end{problem}

\begin{problem}
Write as a percent. Write the remainder in fractional form.
\[
\frac{3}{8}  = \answer{37} \frac{\answer{1}}{\answer{2}}\%
\]
\end{problem}

\begin{problem}
Write as a percent. Write the remainder in fractional form.
\[
1 \frac{5}{9}  = \answer{155} \frac{\answer{5}}{\answer{9}} \%
\]
\end{problem}

\end{document}